% !TEX program = xelatex
\documentclass{英语试卷}
\begin{document}
\examtitle{湖州、衢州、丽水2024年11月三地市教学质量检测试卷}{英语}
\par 湖州、衢州、丽水2024年11月三地市教学质量检测试卷\par
\par 高三年级英语试题\par
\begin{notice}
\par 1. 答卷前，考生务必将自己的姓名、准考证号填写在答题卡上。\par
\par 2. 回答选择题时，选出每小题答案后，用铅笔把答题卡上对应题目的答案标号涂黑。如需改动，用橡皮擦干净后，再选涂其他答案标号。回答非选择题时，将答案写在答题卡上，写在本试卷上无效。\par
\par 3. 考试结束后，将本试卷和答题卡一并交回。\par
\end{notice}
\parttitle{第一部分听力(共两节，满分30分)}
\par 做题时，先将答案标在试卷上，录音内容结束后，你将有两分钟的时间将试卷上的答案转涂到答题卡上。\par
\sectitle{第一节(共5小题：每小题1.5分，满分7.5分)}
\par 听下面5段对话。每段对话后有一个小题，从题中所给的A、B、C三个选项中选出最佳选项。听完每段对话后，你都有10秒钟的时间来回答有关小题和阅读下一小题。每段对话仅读一遍。\par
\question{例：How much is the shirt?}
\begin{choices}(3)
  \item f19.15.
  \item f9.18.
  \item f9.15. 答案是C。
\end{choices}
\question{1. What did Martin do yesterday?}
\begin{choices}(3)
  \item He visited a museum.
  \item He went to a beach.
  \item He had lunch at a café.
\end{choices}
\question{2. How will the man get to Shanghai?}
\begin{choices}(3)
  \item By train.
  \item By plane.
  \item By car.
\end{choices}
\question{3. What does the woman suggest man do?}
\begin{choices}(3)
  \item Avoid eating meat.
  \item Have more beans.
  \item Eat less rice.
\end{choices}
\question{4. What is the probable relationship between the speakers?}
\begin{choices}(3)
  \item Strangers.
  \item Relatives.
  \item Co-workers.
\end{choices}
\question{5. How does the man probably feel now?}
\begin{choices}(3)
  \item Shocked.
  \item Relieved.
  \item Embarrassed.
\end{choices}
\sectitle{第二节(共15小题：每小题1.5分，满分22.5分)}
\par 听下面5段对话或独白。每段对话或独白后有几个小题，从题中所给的A、B、C三个选项中选出最佳选项。听每段对话或独白前，你将有时间阅读各个小题，每小题5秒钟；听完后，各小题将给出5秒钟的作答时间。每段对话或独白读两遍。\par
\question{6. Where did the twins meet?}
\begin{choices}(3)
  \item In Springfield.
  \item In Chicago.
  \item In Seattle.
\end{choices}
\question{7. How much time did Ms.White spend looking for her sister?}
\begin{choices}(3)
  \item Four years.
  \item Forty-five years.
  \item A month.
\end{choices}
\par 听第7段材料，回答第8至10题。\par
\question{8. How many books can a student borrow at most each day?}
\begin{choices}(3)
  \item Ten.
  \item Five.
  \item Three.
\end{choices}
\question{9. Why can't the boy borrow the books now?}
\begin{choices}(3)
  \item He doesn't have his student ID card with him.
  \item He has been fined by the library before.
  \item He has not returned the overdue book.
\end{choices}
\question{10. What does the woman suggest the boy do?}
\begin{choices}(3)
  \item Pay the fine now.
  \item Copy what he needs.
  \item Make a call to his friend.
\end{choices}
\par 听第8段材料，回答第11至13题。\par
\question{11. When did the woman start drawing on stones?}
\begin{choices}(3)
  \item About four years ago.
  \item About four months ago.
  \item When she was four.
\end{choices}
\question{12. What is special about the woman's paintings?}
\begin{choices}(3)
  \item They're about sceneries.
  \item They contain little detail.
  \item Some of them are large.
\end{choices}
\question{13. Why did the woman draw on stones?}
\begin{choices}(3)
  \item To earn more money.
  \item To be famous on the Internet.
  \item To improve her painting skills.
\end{choices}
\par 听第9段材料，回答第14至16题。\par
\question{14. Where was Annie born probably?}
\begin{choices}(3)
  \item In Italy.
  \item In France.
  \item In Portugal.
\end{choices}
\question{15. What is difficult for Annie to learn about Spanish?}
\begin{choices}(3)
  \item Its writing.
  \item Its grammar.
  \item Its pronunciation.
\end{choices}
\question{16. What is Mr.Johnson's opinion of Annie?}
\begin{choices}(3)
  \item She is falling behind.
  \item She is performing average.
  \item She is doing well.
\end{choices}
\par 听第10段材料，回答第17至20题。\par
\question{17. What was the result of the fire?}
\begin{choices}(3)
  \item Many people were injured.
  \item Many factories were destroyed.
  \item Around 10,000 people were homeless
\end{choices}
\question{18. What do people in Mirpur mainly do to make a living?}
\begin{choices}(3)
  \item Make clothing.
  \item Produce plastic.
  \item Sell woods.
\end{choices}
\question{19. What made the fire fighters'work more difficult?}
\begin{choices}(3)
  \item The hot weather.
  \item The lack of water.
  \item The terrible traffic.
\end{choices}
\question{20. Where are the homeless people now?}
\begin{choices}(3)
  \item On the street.
  \item At nearby schools.
  \item In a factory.
\end{choices}
\parttitle{第二部分阅读理解(共两节，满分50分)}
\sectitle{第一节(共15小题：每小题2.5分，满分37.5分)}
\par 阅读下列短文，从每题所给的A、B、C和D四个选项中，选出最佳答案。\par
\passagelabel{A}
\par The International Student Hub(I-Hub)is a centralized and comprehensive center for New York University's international student community.Here are some activities provided by I-Hub for new international students.\par
\par Ikea Trip for New International Students\par
\par Tuesday, August 20th|1:00PM-4:00PM\par
\par Meet up with other new international graduate students to take public transportation to the IKEA in Brooklyn. We will leave from Tandon's Brooklyn campus at 1:00 pm and take the MTA bus to IKEA.Once we arrive at IKEA,you will be free to shop as long as you would like.\par
\par Manhattan Grocery Store Hopping\par
\par Wednesday, August 21st|11:00AM-12:30PM\par
\par Still missing daily essentials but not sure where to buy them? Join us on a trip and visit grocery stores around the Washington Square campus with other new international students! Get insider tips on how to stay under budget when shopping and where to find the best deals.\par
\par Fantastic Night in Coney Island\par
\par Friday, August 23rd |7:00PM-10:00PM\par
\par Looking for a fun and authentic American experience? Join fellow international students at the Brooklyn Cyclones baseball game! It's the perfect way to spend an evening—cheering with passionate audience and soaking up the lively atmosphere at Maimonides Park. There will also be a fireworks show!\par
\par Discover NY: Chinatown\par
\par Sunday, August 25th|11:00AM-1:00PM\par
\par Join us for a walking tour around Chinatown! Come and experience the rich culture of the historic neighborhood and sample a wide selection of cheap eats. On the way, you'll also find out about location attractions, tips on where to buy groceries, and more!\par
\question{21. What can students do in Manhattan Grocery Store Hopping?}
\begin{choices}(2)
  \item Take the MTA bus.
  \item Taste cheap eats.
  \item Watch a fireworks show.
  \item Get advice on purchase.
\end{choices}
\question{22. Which activity will appeal to a sports fan?}
\begin{choices}(2)
  \item Discover NY:Chinatown.
  \item Fantastic Night in Coney Island.
  \item Manhattan Grocery Store Hopping.
  \item Ikea Trip for New International Students.
\end{choices}
\question{23. What do the four activities have in common?}
\begin{choices}(2)
  \item They offer shopping experiences to students.
  \item They require students to take public transport.
  \item They enable students to learn Chinese culture.
  \item They are designed for new international students.
\end{choices}
\passagelabel{B}
\par I switched to a flip phone(翻盖手机). It vastly improved my life. I felt empty and frustrated at how much time I spent on my smartphone. So I spent the last five years steadily stripping my phone down to the basics. Social media was the first to go, then email, then news apps, and then even internet browsers.\par
\par I long considered switching to a flip phone, but I always found an excuse to stay. Ultimately it was two of my students, 15- and 18-year-old sisters, who pushed me over the edge. Word spread around school that they had bought flip phones. If they were brave enough to take the plunge, I could, too. I spent 9 on a flip phone, plus 120 for a year of unlimited talk and text on a new number. For the last five months, I haven't carried a smartphone with me, just my flip phone.\par
\par This is the charming, romantic side of owning a flip phone: more conversations, more eye contact, more connection. The separation between work and life is much more noticeable. I don't overthink work emails all day, since I no longer carry my work email in my pocket. I am living more like my creative, curious, phone-free children.\par
\par This all sounds wonderful because it is wonderful. And now it's confession time: I still have my smartphone. As it turns out, our addiction to smartphones is not just on an individual level. Intentionally or not, we have built structures that depend upon smartphone access, making it an essential feature of modern life. While most of the flip phone inconveniences are charming, some of them are so frustrating that they make you want to pull your hair out.\par
\par Despite that, I will continue to use a flip phone as my primary phone so I don't have to carry a smartphone with me. And I will continue to evangelize for it. I hope many others get flip phones. We will all be better for it. If nothing else, it'll be nice to have more people to talk to in the grocery line.\par
\question{24. What drove the author to switch to a flip phone?}
\begin{choices}(2)
  \item Disturbance of social media.
  \item Overspent time on his smartphone.
  \item Inspiration from his students.
  \item Negative feelings to smartphones.
\end{choices}
\question{25. What happened after the author used a flip phone?}
\begin{choices}(2)
  \item He got less work assignments.
  \item He had more interaction with people.
  \item He spent more time with his children.
  \item He lived a more romantic life.
\end{choices}
\question{26. Why is it difficult to stop using smartphones completely nowadays?}
\begin{choices}(2)
  \item They serve as a vital tool.
  \item They make people overthink.
  \item They represent the modern life.
  \item They bring people many profits.
\end{choices}
\question{27. What is the author's purpose in writing the passage?}
\begin{choices}(2)
  \item To advise readers to use a flip phone.
  \item To advocate living without a smartphone.
  \item To inform readers of a flip phone's advantages.
  \item To stress the disadvantages of having a smartphone.
\end{choices}
\passagelabel{C}
\par The oak tree has been in southern England continuously for 12,000 years. But does it have a future?\par
\par By 2050, London's weather could resemble that of Barcelona, with long stretches of summer drought. These ancient trees were not designed to thrive in such conditions. Around the world, millions of creatures facing unprecedented temperatures and habitat loss are on the move.\par
\par We do not tend to think of trees as migratory: they are slow-growing and rooted to the ground. Faced with this problem, the ecologist Charlie Gardner is one of a growing number proposing a radical, controversial solution: we help the trees on their march. Some scientists advocate moving individual species hundreds of kilometres north, within the country and from abroad. Others say entire ecosystems should be translocated, moving multiple species at once.\par
\par The “ assisted migration" is divisive. Conservationists have long argued against introducing non-native species. They have concerns about unintended consequences of assisted migration.“ Large-scale moving of ecosystems can't be undone(撤销), so should be very much a last choice," said Andrew Allen from the Woodland Trust. Allen believes we should priorities helping native trees adapt to changes by encouraging natural regeneration, expanding and\par
\par connecting existing woodlands.\par
\par But others say those arguments fail to deal with speed or size of change caused by the climate crisis." Earlier in my career the narrative was all about restoring past baselines, and then, as I got more independent in my research, I realised that's impossible. We are changing the climate so much, so we have to be a little more inventive. When we have examples of assisted migration, it's often with species which are absolutely on the edge of extinction and have run out of other options," said Sarah Dalrymple, a conservation ecologist from Liverpool University. who is part of the task force. She wants to see assisted migration explored earlier, working with species before their populations become so small. " You can never get rid of all the risks, but there is also risk in inaction." she said.\par
\par " The risk of inaction is increasing every day."\par
\question{28. What is the problem facing ancient trees in London?}
\begin{choices}(2)
  \item Long summer.
  \item Habitat destruction.
  \item Extreme weather.
  \item Slow growth.
\end{choices}
\question{29. Which of the following statements may Andrew Allen agree with?}
\begin{choices}(2)
  \item Large-scale moving of ecosystems is impossible.
  \item Assisted migration may contribute to bad results.
  \item Non-native species are beneficial to ecosystems.
  \item Native species can survive changes by themselves.
\end{choices}
\question{30. What does Sarah Dalrymple urge people to do in the last paragraph?}
\begin{choices}(2)
  \item Stop changing the climate.
  \item Have more options to help species.
  \item Help native trees to adapt to climate changes.
  \item Take action earlier to prevent species'extinction.
\end{choices}
\question{31. What might be the author's attitude towards assisted migration?}
\begin{choices}(2)
  \item Supportive.
  \item Cautious.
  \item Objective.
  \item Doubtful.
\end{choices}
\passagelabel{D}
\par In a study published in Personality and Social Psychology Bulletin, researchers found that people tend to turn down offers of "free money," as well as unusually high salaries or suspiciously cheap services, because they seem "too good to be true."\par
\par In the initial experiment, nearly 40 percent of participants ate a cookie offered freely - compared with about 20 percent of those offered \$5 as well. Nine further experiments, involving more than 4,000 participants, used online questionnaires to present other scenarios. These included being offered money to accept a ride home, extremely high construction-job wages and surprisingly cheap flights. In each case, past a certain point, higher potential monetary gain reduced participants' likelihood of accepting the offer.\par
\par The study lead author Vonasch says it illustrates that contrary to the "standard economic model," which supposes humans always seek to maximize gains, transactions need to also be understood as social interactions between people trying to understand each other's minds. If someone seems to break accepted rules, such as self-interest, without any explanation, we assume they have hidden motives and infer there will be "phantom costs"\par
\par imagined consequences that reduce what Vonasch calls an offer's" psychological value."\par
\par Factors beyond the present moment may come into play." Understanding that others'perceived overgenerosity may put us in their debt could also help explain people's reluctance," says Rachel McCloy, a psychologist at England's University of Reading. The researchers also showed how to mitigate the effect: simply provide a reason for the deal. The "cheap flights" experiment included a condition where the seats were revealed to be very uncomfortable. "Uncomfortable seats aren't typically a selling point," Vonasch says. "But telling people the seats were uncomfortable made them more willing to take them because it was sufficient explanation."\par
\par The scientists are now experimenting with whether phantom costs play into humans' interactions with robots and artificial intelligence. "If AI is overly generous, will people imagine phantom costs?" Vonasch says. "People tend to treat AI as if they have a mind, when obviously they don't."\par
\question{32. According to paragraph 2, the experiments mainly focus on people's}
\begin{choices}(2)
  \item Thinking method.
  \item Decision making.
  \item Financial condition.
  \item Potential ability.
\end{choices}
\question{33. Why does the author mention "standard economic model"?}
\begin{choices}(2)
  \item To explain the experiments further.
  \item To clarify a psychological concept.
  \item To provide evidence for the study.
  \item To express doubt on the experiments.
\end{choices}
\question{34. What does the underlined word"mitigate" in paragraph 4 probably mean?}
\begin{choices}(2)
  \item Analyse.
  \item Improve.
  \item Measure.
  \item Lessen.
\end{choices}
\question{35. Which of the following is the best title for the text?}
\begin{choices}(2)
  \item Phantom costs:Where" standard economic model"doesn't work
  \item Phantom costs:Whether people's likelihood of accepting an offer reduces
  \item Phantom costs:Why " too good to be true" scares people off
  \item Phantom costs:How imagined consequences prevent hidden motives
\end{choices}
\sectitle{第二节(共5小题；每小题2.5分，满分12.5分)}
\par 阅读下面短文，从短文后的选项中选出能填入空白处的最佳选项。选项中有两项为多余选项。\par
\par Sometimes it can feel like your head is too cluttered. With work, school, or a busy schedule things are bound to get jumbled (乱糟糟的) every once and a while! \blanknum{36} From focusing to journaling, we'll teach you how to mindfully distract yourself from worrying thoughts.\par
\par Focus on the present moment. Practice mindfulness daily to help clear your mind and be more focused. Mindfulness is all about letting all other thoughts come and go. Observe your thoughts and feelings without judgement, as if you're an outsider in your own mind. \blanknum{37}\par
\par Take a walk. \blanknum{38} When you're being mindful, you're fully engaging in the present moment, letting feelings come and go as you focus on what's in front of you. Combining this state of mind with walking produces an extra special stress reliever that improves mental and physical health. Try scheduling a walk around the block on the treadmill, or through a local park every day.\par
\par \blanknum{39} Take a moment to breathe and come back to the present moment. Life can easily become go.go, go, and in the process, you can easily forget to breathe. Settling down wherever you are to take controlled breaths can calm your racing mind and heartbeat, helping you feel more relaxed.\par
\par Journal your thoughts. Writing down your worries can help you sort through them. Try reaching for a notebook or phone when a million thoughts are racing through your mind. \blanknum{40} Journaling can help you be more aware of your emotions and reflect on what's troubling you.\par
\begin{choices}(1)
  \item Do deep breathing exercises.
  \item Take a moment to slow down.
  \item Moving your body mindfully can help things seem clearer.
  \item But you can do something to clear your mind and calm things down.
  \item Set down everything that comes to mind, letting all pour out onto the page.
  \item Then let those thoughts and feelings go and focus on what you're currently doing instead.
  \item Here are some tips on how to protect your head from brain injury and then start your work.
\end{choices}
\parttitle{第三部分语言知识运用(共两节，满分30分)}
\sectitle{第一节完形填空(共15小题；每小题1分，满分15分)}
\par 阅读下面短文，从短文后各题所给的A、B、C和D四个选项中，选出可以填入空白处的最佳选项。\par
\par My story with my cat started many years ago, and now I have to prepare to bid farewell to her. It was 2003 and a childhood 41 was stationed at Fort Bragg. in North Carolina where she was in the 82nd Airborne Division(空降师). She was done jumping out of planes for the day, 42 the base, when she heard meowing coming from a warehouse.\par
\par Inside, she 43 a stray(流浪的)cat, proud but dirty, and decided to take her home. Soon, the cat was 44 . One day, my friend e-mailed me a photograph of the little lovely kittens, arranged on her laps. She 45 , “ Don't you want to have them?"\par
\par A few weeks later, the kittens were 46 onto the back seat of a Jeep and driven north. Two were dropped off at my parents' house, and one with me at my studio apartment. The family 47 was that these were military-grade cats, since their mom was from the base!\par
\par Gradually, I became 48 to my cat. She always leaned on my forearm the moment I told her the 49 of us: I described the apartments we'd lived in, the friends we'd met, and the sunsets we witnessed together. Almost every morning, I 50 with her soft paws rubbing against my face. 51 her bright eyes melted my heart and her shock-white belly in the room was the best 52 throughout my day.\par
\par Twenty-one years later, min e is the 53 one standing. The cat is very sick. so I know a 54 is at hand. Comedy is tragedy plus time. That's the rule. 55 , I think there's a lesson here about" nurture."\par
\cz{41}{friend}{neighbor}{teacher}{classmate}
\cz{42}{setting up}{B.heading off}{C.arriving at}{D.breaking into}
\cz{43}{removed}{B.spotted}{C.identified}{D.adopted}
\cz{44}{restless}{B.injured}{C.pregnant}{D.troublesome}
\cz{45}{added}{B.insisted}{C.declared}{D.explained}
\cz{46}{stuffed}{B.thrown}{C.handed}{D.loaded}
\cz{47}{secret}{B.truth}{C.joke}{D.issue}
\cz{48}{exposed}{B.attached}{C.devoted}{D.limited}
\cz{49}{story}{B.lesson}{C.emotion}{D.theme}
\cz{50}{wept}{B.waited}{C.whispered}{D.woke}
\cz{51}{Instead}{B.Therefore}{C.However}{D.Meanwhile}
\cz{52}{signal}{B.impression}{C.sight}{D.image}
\cz{53}{first}{B.cutest}{C.last}{D.youngest}
\cz{54}{separation}{B.delivery}{C.reunion}{D.celebration}
\cz{55}{Constantly}{B.Personally}{C.Immediately}{D.Obviously}
\sectitle{第二节(共10小题；每小题1.5分，满分15分)}
\par 阅读下面短文，在空白处填入1个适当的单词或括号内单词的正确形式。\par
\par It is considered one of the most extraordinary documentary film \blanknum{56} (project)ever made, a film shot in every country of the world \blanknum{57} the same day, involving 3,000 hours of footage in 70 languages from 19,000 volunteer filmmakers around the world. The 104-minute film, One Day on Earth, is a visual poem \blanknum{58} (star) everyone on the planet. It's about you and me, the times we live in, and our place in \blanknum{59} puzzle of humanity.\par
\par Director Kyle Ruddick and executive producer Brandon Litman \blanknum{60} (be)the two young forces behind One Day on Earth.They met at the University of Southen California,where Ruddick studied film \blanknum{61} Litman majored in business.Although both had worked on short-length \blanknum{62} (commerce)and broadcast projects,they had never attempted anything like One Day on Earth,their first feature-length film,and neither had anyone else.\par
\par Ruddick and Litman came up with the idea for the film in 2008 at a world music festival, \blanknum{63} they heard musicians who had never met before play together for the first time. After a few awkward attempts, the musicians soon discovered a way \blanknum{64} (create)a beautiful fusion of music. In a similar way, "the editorial process was a process of discovery," said Ruddick. "You couldn't make this sort of film without being \blanknum{65} (complete)open to what you receive."\par
\parttitle{第四部分写作(共两节，满分40分)}
\sectitle{第一节(满分15分)}
\par 假定你是李华，你校将举办英语节“戏剧之夜”(Drama Night)活动，并且进行网络直播(live-stream)，请你给英国好友David写一封电子邮件，邀请他观看，内容如下：(1)活动介绍；(2)观看时间和方式。注意：\par
\par (1)写作词数应为80左右；\par
\par (2)请按如下格式在答题纸的相应位置作答。\par
\par\vspace{0.4em}
\begin{center}\begin{tabular}{{|l|}}
\hline
Dear David, Yours, Li Hua \\
\hline
\end{tabular}\end{center}
\vspace{0.4em}
\sectitle{第二节(满分25分)}
\par 阅读下面材料，根据其内容和所给段落开头语续写两段，使之构成一篇完整的短文。\par
\par When Devon's sister, Lil, went off to college, Devon's stomach felt all twisty. Then Mom said that Kapil, a high-school student from India, would be staying in their house and living in Lil's room, for they were the best family. The piece of news made Devon's stomach feel worse, as he missed his sister so much. "What if I can't understand him?" Devon asked. "He speaks English well," Mom said.\par
\par On Kapil's first day, he exchanged some small talk with Devon before disappearing into his new room. Lil's room, Devon reminded himself. As days went on. Kapil started spending more time in the rest of the house. He was in the kitchen a lot, whose cooking always filled the air with steam and spice that made Devon's eyes water. "Smells wonderful!" Mom exclaimed. When Lil called a few days later, Devon told her about Kapil's cooking. "Lucky! I love Indian food!" Lil said. "But it seems too spicy for me," replied Devon.\par
\par One day, Mom had to work late. "Kapil will stay with you until I get home," she explained. Devon's stomach twisted. Lil had always stayed with him when Mom had to work late.\par
\par When Devon got home from school, Kapil greeted him at the door. "Would you like a snack?" Kapil asked. Devon shrugged. Lil and he used to cook together after school. which they called "making messes. Kapil opened the refrigerator." "How about a yogurt?" "I had one at breakfast," responded Devon in a cold voice. Kapil opened a cupboard and smiled. "Wait here." He returned with a package. "I’ll cook an Indian snack," he announced. Kapil pulled out a round, thin cracker the size of a small Frisbee, poured oil into a pan and fried the cracker. Then the kitchen was filled with spicy smells. After a while. Kapil handed Devon a plate with the sizzling pancake on it. "Papadum(印度薄饼)," he said. "Made from bean flour."\par
\par 注意：\par
\par (1)续写词数应为150左右；\par
\par (2)请按如下格式在答题卡的相应位置作答。\par
\writingpara{Paragraph 1: Devon's stomach twisted again.}
\writinglines{8}
\writingpara{Paragraph 2: Looking at the empty plate, Devon said to Kapil, "Let's make messes!" .}
\writinglines{8}
\end{document}
